% META: {"id":"1NSI-TC-CO-018","chapitre":"1NSI-TYPES-CONSTRUITS","type_objet":"corrige","capacites":["1NSI-TYPES-CONSTRUITS-C2"],"mode_creation":"ex_nihilo","sources_inspiration":[],"status":"approved","exercice_ref":"1NSI-TC-EX-018","origin":"needs_review","status_history":["needs_review","audited","accepted","approved"],"release_acceptance":"RELEASE_OWNER_BATCH_ACCEPTANCE","acceptance_closure_digest":"sha256:9b3ccf9a81c5520fbb7e03b7b2d3e7bf2057a02834b6d908bf3be5f49f3b553f","acceptance_receipt_digest":"sha256:4fad86f0282e0fa4e0419cca1ad6379ae7af5a280c04a4a90880f90a1c044242"}
% BEGIN-VERIFY
% def indice_maximum(tab):
%     idx = 0
%     for i in range(1, len(tab)):
%         if tab[i] > tab[idx]:
%             idx = i
%     return idx
% assert indice_maximum([3, 7, 2, 9, 1]) == 3
% assert indice_maximum([10]) == 0
% assert indice_maximum([5, 5, 5]) == 0
% assert indice_maximum([1, 9, 9]) == 1
% END-VERIFY
\begin{corrige}{1NSI-TC-EX-018}
\begin{python}
def indice_maximum(tab):
    # On suppose que l'indice du max est 0
    idx = 0
    # On parcourt les indices a partir de 1
    for i in range(1, len(tab)):
        if tab[i] > tab[idx]:
            idx = i
    return idx
\end{python}

\textbf{Principe :} on initialise \lstinline{idx} à 0, puis on parcourt le tableau à partir de l'indice 1. Si l'élément courant est \emph{strictement} supérieur à \lstinline{tab[idx]}, on met à jour \lstinline{idx}. Le test strict \lstinline{>} (et non \lstinline{>=}) garantit qu'en cas d'égalité, on conserve le plus petit indice.

Vérification : \lstinline{indice_maximum([3, 7, 2, 9, 1])} renvoie \lstinline{3} et \lstinline{indice_maximum([5, 5, 5])} renvoie \lstinline{0}. Résultats confirmés par exécution.
\end{corrige}
